%*******************************************************************************
% Macro per il documento corrente
%*******************************************************************************
\newcommand{\sharedPath}{../shared}
\newcommand{\doctitle}{Business Plan}

%*******************************************************************************
% Preambolo
%*******************************************************************************
\documentclass[a4paper,10pt,twoside]{book}

%*******************************************************************************
% Codifica e lingua
%*******************************************************************************
\usepackage[utf8x]{inputenc}
\usepackage[T1]{fontenc}
\usepackage[english,italian]{babel}
\usepackage{lmodern}
\usepackage{multirow}
\usepackage{amsmath,eurosym}
%*******************************************************************************
% Qualche macro utile a tutti
%*******************************************************************************
\newcommand{\docRoot}{..}
\newcommand{\inglese}[1]{\foreignlanguage{english}{\textit{#1}}}
\newcommand{\team}{\textsf{EtaBeta\,Software}\xspace}
\newcommand{\caName}{BPM-1.0\xspace}
\newcommand{\customer}{\textsf{Alpha\,\&\,Partners}\xspace}
\newcommand{\mktg}{\inglese{marketing}\xspace}
\newcommand{\bsn}{\inglese{business}\xspace}
\newcommand{\sw}{\inglese{software}\xspace}

%*******************************************************************************
% Figure e immagini
%*******************************************************************************
\usepackage{graphicx}
\graphicspath{{\docRoot/shared/pictures/}}

%*******************************************************************************
% Tabelle
%*******************************************************************************
\usepackage{booktabs}
\usepackage{array}

%*******************************************************************************
% Elenchi puntati personalizzati
%*******************************************************************************
\usepackage{enumitem}

%*******************************************************************************
% Landscape per ruotare le pagine
%*******************************************************************************
\usepackage{pdflscape}

%*************************************************
% Collegamenti intra- e intertestuali
%*************************************************
\usepackage{hyperref}
\hypersetup{%
    colorlinks=false,linktocpage=false,pdfborder={0 0 0},%
    pdfstartpage=1, pdfstartview=FitV,plainpages=false,%
    urlcolor=Black, linkcolor=Black,
    pdfcreator={pdfLaTeX},%
    pdfproducer={pdfLaTeX with hyperref package}%
}

%*************************************************
% Grafica vettoriale
%*************************************************
\usepackage{tikz}
\usetikzlibrary{shadows,arrows,shapes}

%*******************************************************************************
% Altri pacchetti
%*******************************************************************************
\usepackage{xspace} % per spazi condizionali extra
\usepackage{lastpage} % per sapere il numero totale di pagine
\usepackage{microtype} % qualche accorgimento tipografico

% *************************************************
% Intestazioni e piè di pagina personalizzati
% *************************************************
\usepackage{fancyhdr}

% stile normale
\fancypagestyle{normal}{
\fancyhead{}
\fancyhead[RE,LO]{
\sffamily\team
}
\fancyhead[LE,RO]{
\sffamily\leftmark
}
\renewcommand{\headrulewidth}{.4pt}
\cfoot{}
\fancyfoot[LE,RO]{\sffamily
  pag. \thepage{} di \pageref{LastPage}}
\fancyfoot[RE,LO]{\sffamily\doctitle}
\renewcommand{\footrulewidth}{.4pt}
}

% stile per gli indici
\fancypagestyle{toc}{
\fancyhead{}
\renewcommand{\headrulewidth}{0pt}
\cfoot{}
\fancyfoot[RO,LE]{\sffamily\thepage{}}
\fancyfoot[RE,LO]{\sffamily\doctitle{}}
\renewcommand{\footrulewidth}{.4pt}
}

\pagestyle{fancy}
\renewcommand{\sectionmark}[1]{\markboth{#1}{#1}}


%*******************************************************************************
% Inizio documento
%*******************************************************************************
\begin{document}

\pagestyle{empty}
\begin{center}

{\sffamily
Sviluppo e Gestione Progetti\\
a.a. 2012--2013
}

\vskip 1.5cm

{
\setlength{\fboxsep}{.2pt}
\fbox{\includegraphics[width=\textwidth]{logo}}
}

\vskip .5cm

{\Huge\sffamily\bfseries
\team
}

\vskip 1.5cm

% titolo del progetto
{\Large\sffamily\bfseries
\caName
}

\vskip 1cm

% titolo del documento
\hrule
\vskip 10pt
{\Huge\scshape
\doctitle
}
\vskip 10pt
\hrule

\end{center}

\clearpage
\pagenumbering{roman} 

\tableofcontents{\thispagestyle{toc}}

\clearpage

\pagestyle{normal}
\pagenumbering{arabic}


% 1) NDA + DISCLAIMER
% 2) Executive summary
% 3) Descrizione generale della nostra azienda  (~3 pagine)
%   3.1) chi siamo
%   3.2) mission/vision
% 4) Mercato di riferimento: 
%   4.1) analisi della domanda (target di riferimento)
%   4.2) analisi del settore (modello delle 5 forze di Porter), SWOT
%   4.3) descrizione del problema + vantaggi competitivi che derivano dalla soluzione
% 5) Descrizione generale del progetto di sviluppo
%   5.1) descrizione di come è articolata soluzione
%       5.2.1) BPM per ottimizzare processi aziendali - ASPETTARE A SCRIVERLO!
%       5.2.2) corso su come utilizzare software BPM
%       5.2.3) acquisizione figure nuove, ad esempio un analista per capire necessità dei clienti e un consulente di psicologia o esperto di psicologia marketing per la formulazione adeguata di proposte
%       5.2.4) corsi di formazione in base alle carenze rilevate (ad esempio web design + accessibilità)
%        5.2.5) corso di Quality Assurance
%   5.2) risorse richieste per la realizzazione del progetto di sviluppo
%       5.3.1) tempistiche
%       5.3.2) risorse umane coinvolte
%       5.3.2) costi
%       5.3.3) prospettive di guadagno negli anni
\chapter{Note preliminari}

\section{Non-Disclosure Agreement}
Nel ricevere questo documento, vi impegnate a mantenere e garantire la massima riservatezza sulle informazioni ivi contenute, e su quelle di cui verrete a conoscenza, anche solo verbalmente, nel corso di eventuali ulteriori indagini e/o incontri, nonché a restituire immediatamente, su richiesta di \customer, tutto il materiale ricevuto senza trattenere alcuna copia.

Questo documento non dovrà essere fotocopiato, riprodotto o distribuito, per intero o in parte, né citato in documenti ufficiali, senza il preventivo consenso scritto di \customer.

\section{Disclaimer}
Il presente \inglese{business plan} è stato redatto secondo ipotesi, dati e indicazioni formulate e fornite da \customer, alla luce delle informazioni note, della situazione in essere e di quanto poteva essere ragionevolmente supposto, al momento della sua stesura.

Si precisa che, in conformità con l'incarico ricevuto, tali informazioni sono state assunte dai materiali redattori acriticamente, ovvero senza svolgere alcun controllo in merito alla correttezza, completezza e validazione dei dati e informazioni ricevute.

\chapter{Executive summary}\label{sec:summary}
% spiegare perché si dovrebbe investire nel progetto, qual è il suo scopo e perché è idealmente "unico"
% evitare slogan

\chapter{Informazioni sulla società}\label{sec:whoweare}

\section{Chi siamo}
\customer è un'azienda che opera nel settore della comunicazione e del marketing sia on-line che off-line. 
L'organizzazione si occupa in particolare di attività di \inglese{copywriter}, \inglese{mktg} tradizionale, grafica, e negli ultimi anni ha aperto le porte anche nei settori della comunicazione visiva e del \mktg on-line e \inglese{social}. In particolare, per quest'ultimo, le prospettive di mercato sono ottime e \customer prevede di aumentare gli investimenti in tale campo.

\customer è nata alla fine degli anni novanta come società di \inglese{marketing} tradizionale. Negli anni ha saputo rinnovarsi e rimanere al passo con un mercato sempre più dinamico e questa è stata la sua arma vincente.

Oggi, pur essendo un'azienda con un numero limitato di risorse, fattura circa  200.000,00 \text{\euro}.

\section{La nostra Mission}

\textit{``Nel mondo l'abito fa il monaco: il nostro obiettivo è farvi l'abito, il monaco lo dovete mettere voi!''}
%Sono riuscita a fuffare solo questo.... idee migliori?!



\chapter{Mercato di riferimento}
%TODO citare sempre fonti!
\section{Analisi della domanda}\label{sec:domanda}
% descrivere mercato target di riferimento e clientela attuale
% dimensioni del target di mercato (in prospettiva diacronica)
% segmentazione del mercato target
% come avviene la distribuzione

\section{Analisi del settore}
% dimensioni e andamento del settore
% barriere all'entrata
% influenze sul settore di cambiamenti macroeconomici
% profittabilità e posizione finanziaria del settore
% ruolo dell'innovazione tecnologica
% influenza di regolamenti/normative
% vantaggio competitivo
% analisi della domanda: trend di crescita/sviluppo della richiesta
% consierare il modello delle 5 forze di Porter (sistema competitivo allargato)
% SWOT analysis

\section{Descrizione del problema}
La \customer , pur essendo un'ottima azienda che riesce, nonostante la crisi attuale del mercato, a fatturare  200.000,00 \text{\euro} all'anno, presenta diversi problemi relativi allo svolgimento dei processi.
I processi costituiscono il cuore di un'azienda e quindi, per mantenere l'azienda un \inglese{leader} del mercato in cui opera, è fondamentale avere una buona organizzazione dei processi.
I problemi che si sono rilevati all'interno dei processi di \customer si riferiscono in particolare alle seguenti fasi nel ciclo di realizzazione di un progetto di consulenza:
\begin{itemize}
	\item Analisi delle necessità;
	\item Predisposizione offerta;
	\item Negoziazione e sottoscrizione offerta.
\end{itemize}

L'organizzazione ha rilevato la presenza di diverse problematiche relative alla fase di analisi di necessità riguardanti le richieste di un cliente. Vi sono stati infatti, diversi casi in cui, \customer  ha dovuto eseguire delle modifiche ai progetti di consulenza a sue spese, in quanto, non avendo compreso a pieno le necessità del cliente, ha  realizzato in maniera inadeguata la consulenza. 

Sono inoltre stati incontrati dei problemi anche nella stesura delle proposte di consulenza. La forma con cui si predispone un'offerta è molto importante. Pur essendo il contenuto di una proposta la fonte primaria di accettazione da parte del cliente, la forma con cui si espone il contenuto influenza il verdetto finale.
Da questo derivano inoltre problematiche relative alla negoziazione e sottoscrizione di un'offerta. La \customer , spesso, non è in grado di ``condizionare'' il cliente e portare la negoziazione ad una convergenza vantaggiosa per entrambe le parti.

Il piano di sviluppo avrà quindi l'obiettivo di risolvere tali problematiche, aumentare l'efficienza e la produttività in modo tale da permettere ad \customer di ricoprire un ruolo da protagonista nello scenario di mercato degli anni a venire.

%TODO questa è la parte da integrare in qualche modo
Nel corso dell'ultimo quinquennio, complice la situazione di criticità generalizzata venutasi a creare nei settori in cui operano i principali clienti della società descritti nella sezione~\ref{sec:domanda}, il modello di \bsn di \customer ha avuto modo di dimostrare una serie di limiti strutturali.

Una condizione dominata dalla costante incertezza rende infatti impossibile elaborare strategie di lungo periodo anche nell'evasione delle commesse, dal momento che a fronte di un volume costante di richieste di quotazioni da parte dei clienti si è registrata una flessione delle richieste di prestazione, in altre parole, sempre più spesso giungono richieste di preventivi ma il conseguente lavoro non viene confermato facendo sì che il tempo e l'impegno investiti nella preventivazione e nella valutazione di fattibilità si rivelino una fonte di perdita economica anziché di guadagno.

 % vantaggi competitivi che derivano dalla soluzione

\chapter {Descrizione generale del progetto di sviluppo}
\section{Descrizione di come è articolata soluzione}
\subsection{BPM per ottimizzare processi aziendali}

\subsection{Corso su come utilizzare software BPM}

\subsection{Acquisizione figure nuove}
%, ad esempio un analista per capire necessità dei clienti

\subsection{Corsi di formazione in base alle carenze rilevate} 
%(ad esempio web design + accessibilità)

\subsection{Corso di Quality Assurance}
  
\section{Risorse richieste per la realizzazione del progetto di sviluppo}
\subsection{Tempistiche}

\subsection{Risorse umane coinvolte}

\subsection{Costi}

\subsection{Prospettive di guadagno negli anni}

\end{document}
